\section{Background and Motivation}
\label{sec:motivation}

\subsection{Why Real Peripherals Matter}

Full-system validation exercises more than a register specification. An NVMe
driver allocates submission and completion queues, constructs PRP lists, rings a
doorbell, waits for DMA and then consumes a phase-tagged completion. A NIC driver
programs descriptor rings, transfers ownership with queue-tail registers, and
interprets device-specific interrupt causes. Controller firmware, reset corner
cases, DMA ordering, and the interaction with an unmodified driver all affect
the observed behavior. A software model is useful for early bring-up and fault
injection, but every modeled behavior is also an assumption. A physical backend
provides an independent implementation of the device data plane and firmware.

\subsection{Why Physical Attachment Is Insufficient}

Conventional attachment couples three decisions that validation engineers need
to vary independently:
\begin{enumerate}[leftmargin=*]
    \item \textbf{Presentation:} which topology, identities, BARs, and
    capabilities the DUT discovers;
    \item \textbf{Control realization:} where configuration and
    protocol-specific control semantics execute;
    \item \textbf{Data realization:} how a device reaches DUT memory and how
    data and completions return.
\end{enumerate}
The board topology commonly fixes all three. PCIe extension and disaggregation
can move the physical device, but do not automatically synthesize the DUT-facing
topology or translate guest DMA addresses and device semantics. SCOPE therefore
defines the software-configurable boundary narrowly: the FPGA bitstream provides
the generic ECAM, route, mailbox, DMA-window, and interrupt mechanisms, while a
startup backend configuration selects the active endpoints and their bindings.
This is configuration-only composition within a pre-provisioned hierarchy, not
arbitrary runtime hot-plug or an unlimited number of physical devices.

\subsection{The Peripheral Integration Tax}

We define the \tax as the incremental hardware and engineering cost of adding a
device type or instance to an FPGA validation platform. It includes device-
specific RTL, configuration state, on-chip buffers, routing pressure, timing
closure risk, and resynthesis effort. Two axes matter. A new \emph{type} may
require new protocol state machines, whereas another \emph{instance} requires
more addressable state and routing capacity. No implementation can make the
second cost zero. Our goal is narrower: keep the FPGA-side per-instance cost to
generic state and routing entries, and move device-specific state growth to the
software back-end.

\subsection{Requirements}

Four requirements follow. First, an unmodified DUT OS must enumerate and bind
its native driver. Second, configuration, BAR routing, DMA, completion, and
interrupt state must remain mutually consistent even though different agents
own them. Third, control mediation must not force all payload bytes through the
same software path. Fourth, adding an endpoint within the provisioned capacity
should not require a different DUT interface or bitstream. The last requirement
is an implementation property of the generic mechanisms; it still needs a
capacity sweep before we claim a particular FPGA resource slope. \sys does not
attempt to reproduce a physical PCIe link: it validates the system-software and
device interaction above the link layer, not PHY, LTSSM, DLLP, or electrical
behavior.
